\section{A geometry that stays fixed during iteration}
\label{sec:fixed-geometry}
\label{sec:rosen-physical-grading}
\label{sec:optical-graph-height}
\label{sec:pp-positive-exact-poisson}

The reference inverse is useful only if a nearby equation can be compared
with it. Two features would normally move: the circles where the equation
changes type, and the angular positions at which characteristics arrive.
Section~\ref{sec:scalar-core} fixes the first by a correction proportional
to the curvature residual. We shall fix the second in the same way.
The resulting coordinates will put all the operators on one set of
cylinders with one set of transmission conditions.

This construction involves infinitely many components. Before moving
coordinates, we explain the quantitative room available on them. The
point is that the deformation dies out much faster than the reference
slope. This room pays for every fixed number of coordinate derivatives.

\subsection{The three scales and the function spaces}

Let $d_\partial$ be distance to the nearest endpoint. With $d_0>0$ fixed,
put, sufficiently near that endpoint,
\begin{equation}\label{eq:readable-physical-weights}
 x_\partial=\sqrt{\log(d_0/d_\partial)},\quad
 H=e^{-2x_\partial^3/3},\quad
 E_0=e^{1+x_\partial^{5/2}\log(e+x_\partial)},\quad
 M=e^{1+x_\partial^{5/2}\log^2(e+x_\partial)}.
\end{equation}
Use smooth positive extensions through the middle. The endpoint
construction gives $x_\partial=n+O(1)$ on the $n$th component. Thus
$H$ agrees with $H_n$, and $E_0$ with $\mathcal E_n$, up to powers of
$\mathcal E_n$. The slope itself has the slower decay
$w_*=\exp(-n^{5/2}+O(n^2))$.

There are two useful inequalities, for every fixed $A$ and $\epsilon>0$:
\begin{equation}\label{eq:full-material-envelope}
 E_0^A\le C_A M,\qquad H^\epsilon M^A\le C_{A,\epsilon}.
\end{equation}
They follow by comparing logarithms. The additional logarithm in $M$
is deliberate: however many fixed coefficients appear after
differentiation, their \emph{whole product} costs one $M$ order.
That order does not grow with the number of derivatives. The cubic
exponent in $H$ still dominates every fixed power of $M$.

Choose
\[
 \ell=\frac{d_\partial}{1+x_\partial},\qquad
 d\varrho=\frac{dr}{\ell},\qquad \mathsf D=\partial_\varrho=\ell\partial_r
\]
near the ends, with a positive extension in between. Then
$\varrho$ covers $\R$ and
\[
 |\mathsf D x_\partial|=\frac1{2x_\partial(1+x_\partial)},\qquad
 |\mathsf D\log H|=\frac{x_\partial}{1+x_\partial}.
\]
All positive $\mathsf D$ derivatives of $\log H,\log M,\log\ell$ are
bounded. Ordinary radial differentiation costs factors bounded by
$E_0^{C_s}$.

We take $\beta=3/4$, within the convenient interval
\begin{equation}\label{eq:beta-choice}
 \frac12<\beta<1.
\end{equation}
The upper inequality leaves two neighboring exponents for the reference
inverse. The lower one makes a quadratic correction decay faster than
$H$, giving the strengthened metric estimate in
Theorem~\ref{thm:new-annular-height}.
The words weak and strong below compare endpoint weights: both
scales consist of smooth functions with flat endpoint extensions.

\begin{definition}[Height, source and circle spaces]
\label{def:readable-physical-grading}
For real even $\pi$-periodic functions set
\begin{equation}\label{eq:full-physical-norm}
 \|u\|_{\gamma,s}^2
 =\sum_{a+b+c\le s}\|M^c\mathsf D^a\partial_\theta^b
                  (H^{-\gamma}u)\|_{L^2(d\varrho\,d\theta)}^2.
\end{equation}
The intersection over $s\ge0$ is $\mathcal X^0_\gamma$.
Choose $\chi_L$ equal to one at the left end and zero at the right,
with derivative supported in a compact positive component, and set
$\mathcal X_\gamma=\mathcal X^0_\gamma\oplus\R\chi_L$.
Its extra scalar records the difference of the two endpoint heights.
The compatible source space is
\[
 \mathcal Y_\beta=\{f\in\mathcal X^0_\beta:
       \mathsf Qf=0\},\qquad
 \mathsf Qf=\int_{r_-}^{r_+}\av f\,dr.
\]
For circle sequences define
\begin{equation}\label{eq:weak-circle-spaces}
 \|a\|_{\mathcal Z_{\gamma,s}}^2
 =\sum_m\sum_{b+c\le s}M_m^{2c}H_m^{-2\gamma}
                    \|\partial_\theta^b a_m\|_{L^2(\T)}^2.
\end{equation}
The angularly constant sector is denoted by $\mathcal K_\gamma$.
Odd and even sectors will be indicated explicitly. Without a displayed
$\gamma$, a height or source norm uses $\gamma=\beta$.
\end{definition}

\begin{lemma}[Calculus and flat extension]
\label{prop:full-physical-tame}
These spaces are tame Fr\'echet spaces. Products, finite-order
differential expressions, and coefficient reciprocals after separation
of their reference factors are smooth tame on neighborhoods where the
remaining denominators have positive lower bounds. The loss is fixed
independently of the output order. Every member of
$\mathcal X^0_\gamma$, $\gamma>0$, extends smoothly and flatly in $r$
at both endpoints; this also holds for its parameter derivatives.
\end{lemma}

\begin{proof}
Here the noncompact weighted scale, rather than the local Sobolev
calculus, needs explanation. Write $v=H^{-\gamma}u$. Choose unit-scale
functions $\phi_i(\varrho)$ with bounded derivatives, finite overlap
and $\sum_i\phi_i^2=1$. Bounded logarithmic derivatives imply
$M\asymp M_i:=M(i)$ on their supports. The norm of $v$ is equivalent to
\[
 \sum_{i,j}\int (M_i^2+\xi^2+j^2)^s
                    |\widehat{\phi_i v}(\xi,j)|^2\,d\xi.
\]
Let $\omega_i=(M_i^2+\xi^2+j^2)^{1/2}$. A smooth partition in
$\log\omega_i$ gives maps $J,P$ with $PJ=I$ to and from the
exponentially decreasing sequence space over
$B=\ell_i^2L_\xi^2\ell_j^2$. Concretely,
$(Jv)_{k,i}=\zeta_k(\omega_i)\widehat{\phi_i v}$ and
$Pb=\sum_i\phi_i\mathcal F^{-1}\sum_k\widetilde\zeta_k(\omega_i)b_{k,i}$,
where the enlarged cutoff is one on the support of $\zeta_k$.
Cauchy--Schwarz in $k$ gives
$\|Jv\|_{s}\le C_s\|v\|_{s+1}$; finite overlap gives
$\|Pb\|_s\le C_s\|b\|_s$. This is a tame retraction and supplies
smoothing. Parity projection preserves it. The endpoint scalar is a
finite-dimensional summand. If $\psi$ is compactly supported with
$\mathsf Q\psi=1$, then $f\mapsto f-\psi\mathsf Qf$ gives the source
projection: its continuity follows by pairing $H^{-\beta}f$ with the
square-integrable fixed function $H^\beta\ell$.

On each unit cylinder the usual Sobolev product bound, at a fixed
base order $b>1$, puts only one factor at the high order. Summing with
finite overlap gives the same bound globally. Derivatives of a
reciprocal are products divided by a positive power of its denominator;
interpolation puts intermediate derivatives into the same one-high-factor
form. After expanding a complete expression, all its reference
coefficients have a bound $E_0^{C_s}$; the first inequality in
\eqref{eq:full-material-envelope} costs one additional weight order.
The differential expression has fixed order, so the total loss is fixed.
The identical Fourier construction proves the assertion for circle
sequences. A restriction to a transverse circle costs one normal
Sobolev derivative, and a fixed cutoff supplies its right inverse.

Finally $\partial_r^a$ is a finite combination of $\mathsf D$ derivatives
with coefficients bounded by $E_0^{C_a}$. Local Sobolev embedding bounds
$d_\partial^{-k}\partial_r^a\partial_\theta^b u$ by a seminorm times
$H^\gamma E_0^{C_{a,b,k}}$, which tends to zero. This proves the flat
extension, also after any fixed parameter differentiation.
\end{proof}

\subsection{Quantitative root geometry}

Write $\mathfrak r_n$ for the radius scale of a model, bounded above
and below, and $v=r_X$. The radial construction gives
\[
 A_n\asymp e^{-n^2},\quad B_n\asymp A_n^9,\quad
 c\mathfrak r_n A_n^5\le v\le C\mathfrak r_n A_n,\quad
 w_* =\lambda_n v,\quad \log\lambda_n=-n^{5/2}+O(n).
\]
The long positive portion and the negative portion have radial size
comparable to $A_n$. Summing the remaining lengths gives
$cA_{n+1}\le d_\partial\le CA_n$, which proves the comparison used
above. On a bounded enlargement the index changes by $O(1)$, so the
same comparisons hold. Every interval length and reciprocal is bounded
by $e^{Cn^2}$.

The derivative and reciprocal bounds were proved in
\eqref{eq:geometry-profile-jets}. They imply that every fixed jet of
$w_*^{\pm1}$ costs only $\mathcal E_n^{C_s}$. We shall also use the
model identities $r_X=v$, $v_X=-\rho r$, with
$\rho=A_n(e^{-X}-e^{-2X})-B_n$, to control the roots quantitatively.

Let $r_{n,L},r_{n,R}$ be the roots of $w_*'$ and let $x$ measure
distance into a negative component. Its reference characteristic speed is
\[
 v_0=\sqrt{-w_*'/(rw_*)}=\sqrt{x}\,c_i(x).
\]
At a root,
$c_i(0)^2=|\rho_X(X_i)|/v(X_i)^3$. Since the left and right values
of $|\rho_X|$ are respectively comparable to $A_n,A_n^9$, while
$v$ there is comparable to $\mathfrak r_n A_n,
\mathfrak r_n A_n^5$, we obtain
\begin{equation}\label{eq:geometry-root-scales}
 c_L(0)\asymp\mathfrak r_n^{-3/2}A_n^{-1},\qquad
 c_R(0)\asymp\mathfrak r_n^{-3/2}A_n^{-3}.
\end{equation}
Dividing a simple zero by its signed distance uses
$f(x)/x=\int_0^1f'(tx)\,dt$. It follows from the model equations that
the normalized jets of $c_i$ and $c_i^{-1}$ on a collar of size
$c\mathfrak r_n A_n^5$ cost at most
$C_k(\mathfrak r_n A_n^5)^{-k}$; at the left root $A_n$ suffices.
Choose once and for all
\[
 \delta_n^{\mathrm{root}}=\epsilon\mathfrak r_n A_n^{20}.
\]
On $|x|<8\delta_n^{\mathrm{root}}$ the ratio $c_i(x)/c_i(0)$ lies in $[1/2,2]$.
Higher derivatives change constants, not the collar. The fixed central
pieces are handled by their positive minima.

\begin{lemma}[Travel-angle controls]
\label{lem:quantitative-optical-controls}
There are profiles $\chi_n$ supported in $\delta_n^{\mathrm{root}}<x<2\delta_n^{\mathrm{root}}$
at the left root, with
\[
 DT_n[\chi_n]=1,
 \qquad \|\partial_r^k\chi_n\|_\infty
       \le C_k\mathfrak r_n^{-k}A_n^{9-20k},
 \qquad T_n(w)=\int_{w'<0}\sqrt{-w'/(rw)}\,dr.
\]
The differential uses the variation $w_*e^{t\chi_n}$.
A separate correction cutoff can rise from zero to one on
$3\delta_n^{\mathrm{root}}<x<4\delta_n^{\mathrm{root}}$. On both supports,
$v_0\asymp\mathfrak r_n^{-1}A_n^9$ and
$|w_*'/w_*|\asymp\mathfrak r_n^{-1}A_n^{18}$.
\end{lemma}

\begin{proof}
For nonnegative nonzero $\phi\in C_c^\infty((1,2))$ take
$f_n=\phi(x/\delta_n^{\mathrm{root}})$. Logarithmic variation gives
$\delta v_0=-f_n'/(2rv_0)$, hence
\[
 DT_n[f_n]=\frac12\int((rv_0)^{-1})'f_n\,dr
 =-\frac{(\delta_n^{\mathrm{root}})^{-1/2}}{4r_{n,L}c_L(0)}
       \int_1^2t^{-3/2}\phi(t)\,dt\,(1+O(A_n^{19})).
\]
This is nonzero and comparable to $-A_n^{-9}$.
Set $\chi_n=f_n/DT_n[f_n]$. The derivative bound follows by rescaling.
The speed bounds follow from \eqref{eq:geometry-root-scales}; the
second is $|w_*'/w_*|=rv_0^2$. The disjoint cutoff is an integrated
rescaled bump. The same argument at each fixed central root completes
the finite part.
\end{proof}

\subsection{One neighborhood for all components}

The unknown metric has the form
\begin{equation}\label{eq:full-height-metric-controls}
 w=w_*\exp\left(\nu\chi_*+\sum_m\kappa_m\chi_m\right).
\end{equation}
Here $\chi_*$ is the compact positive control of the second mode.
The remaining controls have disjoint negative supports. Put
$\Phi_*=\phi_*\sqrt2\cos2\theta$, with $\ell_*\Phi_*=1$.
Our state is
\begin{equation}\label{eq:rewrite-state}
 \mathbf x=(u,\nu,\kappa,h_L),\quad
 u\in\Phi_*+\ker\ell_*\subset\mathcal X_\beta,\quad
 \kappa\in\mathcal K_\beta,\quad
 h_L-\id\in\mathcal Z_{\beta,\mathrm{odd}}.
\end{equation}
Each chart is an increasing lift commuting with translation by $\pi$.
The full displacement seminorm is
\begin{equation}\label{def:joint-displacement-seminorms}
 \mathfrak n_s=|q|+\|u-\Phi_*\|_s+|\nu|
           +\|\kappa\|_{\mathcal K_{\beta,s}}
           +\|h_L-\id\|_{\mathcal Z_{\beta,s}}.
\end{equation}
The unscaled $u$ must appear here even at $q=0$, since the metric
derivative of $\mathcal F_0(w,u)=\mathscr L_wu$ depends on $u$.

For local norms include the component and all neighboring components
used by its charts; for a positive component include both adjacent
negative ones. These neighborhoods have bounded overlap. Compact
scalars are included only where their supports meet a neighborhood.
Let $\mathfrak g_{s,m}$ be the corresponding geometric norm of
$(qu,\nu,\kappa,h_L-\id)$, using only positive height jets at an end.
Then $(\sum_m\mathfrak g_{s,m}^2)^{1/2}\le C_s\mathfrak n_s$.

\begin{proposition}[Uniform geometric calculus]
\label{prop:uniform-physical-geometry}
\label{lem:common-geometric-neighborhood}
\label{lem:current-geometric-jets}
On a single neighborhood $\mathfrak n_{b_0}<\epsilon_0$, for fixed
finite $b_0$, all metric type intervals are unchanged. With
$Z=qu$, $E=1+w^2$,
\[
 Q=w+Z_r+EZ_{\theta\theta}/r,\qquad
 D=1-2wZ_r-Z_r^2-E(Z_\theta/r)^2,
\]
one has $w/w_*\in[1/2,2]$, $Q/w\in[7/8,9/8]$, $D\in[1/2,3/2]$
and $1-(Z_\theta/r)^2\ge3/4$.
The regularized coordinates, inverses, height and source factors
constructed below have reference jets bounded by
$C_s\mathcal E_m^{C_s}$. Their differences satisfy
\begin{equation}\label{eq:geometry-current-jet-bound}
 \|\Delta c_m\|_{C^s}
 \le C_s\mathcal E_m^{C_s}H_m^\beta\mathfrak g_{s+d_*,m}
\end{equation}
with fixed $d_*$. The assertion includes the positive principal
differences divided by $y^3(1-y)^3$ and every fixed number of state
variations, with the corresponding one-high-factor tame estimates.
\end{proposition}

\begin{proof}
We prove the low comparisons now and the coordinate assertion as its
formulas are constructed below. Local Sobolev embedding and
\eqref{eq:geometry-profile-jets} bound every required low positive
height jet, control jet and chart displacement by
$C\epsilon_0H_m^\beta\mathcal E_m^C$. For each fixed $A$,
$H_m^\beta\mathcal E_m^A$ is bounded and tends to zero. Thus one
$\epsilon_0$ suffices for every component and every inequality in a
finite list. On a control support the relative change of $w'/w$ is
at most $C|\kappa_m|A_n^{-29}$ by the preceding lemma; near each
root it is zero. The compact control has a positive reference minimum.
This preserves the signs and all root germs. The formulas for $Q,D$
and the reciprocal bound for $w_*$ give the remaining comparisons.

Here is the uniform derivative argument used for the forthcoming
coordinate constructions. In a characteristic equation
$\dot\Theta=G(r,\Theta)$ the angular Jacobian is
$\exp\int G_\theta(r,\Theta)\,dr$. The reference has $G_\theta=0$;
the perturbed integral is bounded by
$C\epsilon_0H_m^\beta\mathcal E_m^C$. It therefore stays uniformly
bounded above and below. Highest flow derivatives satisfy a linear
equation with this same integrating factor. Subtracting the reference
equation places a difference in each source term. Product estimates
give \eqref{eq:geometry-current-jet-bound}, with one high factor.
Inversion follows by differentiating $\Psi^{-1}\Psi=\id$, since
the low differential is relatively close to the reference.
Products, positive powers, the finite boundary jets and the divided
differences below have a fixed derivative cost. They preserve the
difference factor $H_m^\beta$. At the end of each calculation absorb
the full $\mathcal E_m^{C_s}$ product into one $M$ order.

All parameter variations use the same formulas. Pulling back by
$Z=qu$ uses $DZ=\dot q\,u+q\dot u$ and
$D^2Z=q_1u_2+q_2u_1$, with higher derivatives zero; thus the full
seminorms \eqref{def:joint-displacement-seminorms} apply also at $q=0$.
The explicit nondegeneracy and finite derivative counts needed in this
argument are verified below, completing the proposition without an
additional smallness condition at higher orders.
\end{proof}

\subsection{Fixing characteristic returns by their own residual}

Start with the type-preserving operator $\bar L$ of
Lemma~\ref{lem:type-completion}. On a negative component write its
coefficients as $(a,b,c)$ and put
\[
 d=b/a,\qquad v=\sqrt{b^2-ac}/a,\qquad c_\pm=d\pm v.
\]
Reflection $S\theta=-\theta$ gives $c_+=-c_-\circ S$.
Let $f$ be the left-to-right minus flow and $T$ its reference travel
angle. Write $R_T\theta=\theta+T$. Define
\begin{equation}\label{eq:full-height-one-chart}
 \widetilde h_R=R_{-T}h_Lf^{-1},\quad
 h_R=\mathsf O\widetilde h_R,\quad
 \mathfrak e=(\mathsf E\widetilde h_R)\circ h_R^{-1},
\end{equation}
where $\mathsf O$ and $\mathsf E$ are odd and even projections.
The derivative of $h_R$ is an average of positive derivatives, so
$h_R$ is increasing. Direct multiplication gives
\[
 (\id+\mathfrak e)h_Rfh_L^{-1}=R_{-T}.
\]
Thus one even equation $\mathfrak e=0$ imposes both returns, the plus
one by reflection. Its mean controls the travel angle; its nonconstant
part controls the boundary parametrization.

For an odd variation $\delta h_L=\psi\circ h_L$, at fixed $f$, the
chain rule yields
\begin{equation}\label{eq:full-height-chart-derivative}
 \begin{split}
 D_h\mathfrak e[\psi]&=\mathcal M_T\psi+\mathcal Q_h(\mathfrak e,\psi),\qquad
 \mathcal M_T\psi=\mathsf E[\psi(\theta+T)],\\
 \mathcal Q_h(\mathfrak e,\psi)&=
 \mathsf E\left[\mathfrak e\int_0^1
       \psi'(\theta+T+t\mathfrak e)\,dt\right]
 -\mathfrak e'\mathsf O[\psi(\theta+T+\mathfrak e)].
 \end{split}
\end{equation}
Indeed the variation of $\widetilde h_R$, after composition with
$h_R^{-1}$, is $\psi(\theta+T+\mathfrak e)$; its odd part is the
variation of $h_R$ in these coordinates. This proves the formula.
At the reference,
\begin{equation}\label{eq:rewrite-circle-inverse}
 \mathcal M_T^{-1}(g_j\cos j\theta)
       =\frac{g_j}{\sin(jT)}\sin j\theta\quad(j>0\text{ even}).
\end{equation}
The selected angle bound loses one angular derivative and a fixed
coefficient weight. There is no inverse on the constant mode; this
is why the radial controls were added.

\begin{proposition}[Return correction]
\label{prop:optical-residual-completion}
The operator $\bar L$ can be replaced, on the chosen neighborhood,
by a symmetric divergence operator $\widehat L$ with the same normal
coefficient, the same root germs, and exact returns
$h_R\widehat f_\pm h_L^{-1}=R_{\pm T}$. It is unchanged on positive
components. Its difference has the form
\begin{equation}\label{eq:optical-residual-completion}
 \bar L-\widehat L=\mathscr E_{\mathbf x}(\mathfrak e),
\end{equation}
where $\mathscr E_{\mathbf x}(e)\eta$ is smooth tame and bilinear
in $(e,\eta)$, with values in $\mathcal Y_\beta$.
\end{proposition}

\begin{proof}
Let $\Phi_r$ be the minus flow. Its required endpoint is
$f_*=h_R^{-1}R_{-T}h_L$. With $H_c=h_Rf$, set
\[
 K=f^{-1}f_*=H_c^{-1}(\id+\mathfrak e)H_c,
 \quad\Delta=K-\id,\quad J_r=\id+\chi(r)\Delta,
 \quad\widehat\Phi_r=\Phi_rJ_r.
\]
Use the interior cutoff of the control lemma. Each $J_r$ is increasing,
as a convex combination of increasing lifts. The new velocity is
\begin{equation}\label{eq:readable-optical-velocity}
 \widehat c_-=c_-+
 \Phi_r'(\Phi_r^{-1}\theta)\chi'(r)
                   \Delta(J_r^{-1}\Phi_r^{-1}\theta).
\end{equation}
Define $\widehat c_+=-\widehat c_-\circ S$, then
$\widehat d=(\widehat c_++\widehat c_-)/2$,
$\widehat v=(\widehat c_+-\widehat c_-)/2$,
$\widehat a=a$, $\widehat b=a\widehat d$ and
$\widehat c=a(\widehat d^2-\widehat v^2)$.
These formulas give the asserted returns and symmetric principal matrix.
The speed lower bound on the correction support and the factor
$H_m^\beta\mathcal E_m^C$ make the correction less than $v/2$
on a single low-order neighborhood. Near the roots it is zero.

For the error factorization use
\[
 \Delta(\theta)=\mathfrak e(H_c\theta)
 \int_0^1(H_c^{-1})'(H_c\theta+t\mathfrak e(H_c\theta))\,dt.
\]
Freeze the state and these composition paths. Replace only the leading
residual by an independent $e$. This defines linear increments
$d_1[e],v_1[e]$ with actual values $d_1,v_1$ at $e=\mathfrak e$.
Then the coefficients of $\bar L-\widehat L$ are
\[
 \delta b[e]=-a d_1[e],\qquad
 \delta c[e]=-a(2d+d_1)d_1[e]+a(2v+v_1)v_1[e].
\]
The associated divergence expression is bilinear in $(e,\eta)$.
Its support is interior to each component and all its endpoint jets
are flat. Its integral is zero. Each term contains a residual factor
and a positive height jet, both of weight $H_m^\beta$; dividing by
the source weight leaves enough decay for every fixed coefficient
product. The flow estimates above and the product calculus prove the
claimed smooth tame bounds.
\end{proof}

Combining the two residual corrections gives the exact identity
\begin{equation}\label{eq:readable-completed-height-derivative}
 D_u\mathcal F_q\eta=\widehat L\eta
       +q\partial_\theta(B\mathcal A\eta_\theta)
       +\mathscr E_{\mathbf x}(\mathfrak e)\eta,
 \qquad \mathcal A=\mathcal F_q(w,u),\quad B=E/(rQ).
\end{equation}
Both added terms vanish at a solution. We shall invert $\widehat L$
while retaining these errors in the nonlinear theorem.

\subsection{The root coordinate and the fixed transmission law}

Let $S_\pm$ be inverse characteristics of $\widehat L$, equal to
$\theta$ at the left boundary. Set
\begin{equation}\label{eq:fixed-strip-coordinates}
 U=h_LS_+,\quad V=h_LS_-,\quad
 \sigma=(V-U)/2,\quad\varphi=(U+V)/2.
\end{equation}
Their endpoint values are $(0,h_L)$ and $(T,h_R)$. Moreover
$U_\theta,V_\theta>0$ and
\[
 \det D(\sigma,\varphi)=vU_\theta V_\theta>0.
\]
At fixed $\varphi$, $\sigma$ increases strictly from $0$ to $T$;
this proves global invertibility onto the cylinder. Substitution in
the principal quadratic form, followed by division by this Jacobian,
gives the principal divergence density
$\delta\operatorname{diag}(1,-1)$, where $\delta=\sqrt{b^2-ac}$.

At a root $v=\sqrt x\,c_0(x,\theta)$. The travel coordinate is
of order $x^{3/2}$, so use
$\xi_L=(3\sigma/2)^{2/3}$ and
$\xi_R=(3(T-\sigma)/2)^{2/3}$ on the two halves. To justify
smoothness when the coefficients depend on angle, solve the fixed-time
equation in independent variables $x,\lambda$:
\begin{equation}\label{eq:fixed-strip-root-flow}
 \Theta_t=2xt\,d(xt^2,\Theta)+2\lambda t^2c_0(xt^2,\Theta),
 \qquad 0\le t\le1.
\end{equation}
The inverse time-one flow followed by the boundary chart is one smooth
map $G(x,\lambda,\theta)$. The two characteristic families are obtained
from this same map by $\lambda=\pm x^{3/2}$. Their common boundary
chart is essential for the cancellation in the following estimate.

\begin{lemma}[Finite-order control at a root]
\label{lem:rewrite-root-composition}
On a fixed rectangle $0\le x\le1$, $|\lambda|\le1$, define, for a
smooth $G(x,\lambda,\theta)$,
\[
 \mathcal E G=\frac{G(x,x^{3/2},\theta)+G(x,-x^{3/2},\theta)}2,\qquad
 \mathcal O G=\frac{G(x,x^{3/2},\theta)-G(x,-x^{3/2},\theta)}
                         {2x^{3/2}}.
\]
Both extend smoothly to $x=0$, and for every integer $s\ge0$,
\begin{equation}\label{eq:rewrite-root-composition-bound}
 \|\mathcal E G\|_{C^s_{x,\theta}}\le C_s\|G\|_{C^s_{x,\lambda,\theta}},
 \qquad
 \|\mathcal O G\|_{C^s_{x,\theta}}
                       \le C_s\|G\|_{C^{s+1}_{x,\lambda,\theta}}.
\end{equation}
The quotient at zero is $\partial_\lambda G(0,0,\theta)$.
The estimates also hold after any fixed number of parameter
differentiations.
\end{lemma}

\begin{proof}
First suppose $G$ is even in $\lambda$. Fix $0<x_0\le1$ and
take its degree-$s$ Taylor polynomial at $(x_0,0)$.
Its monomials $(x-x_0)^i\lambda^{2j}$ become
$(x-x_0)^ix^{3j}$. Their $a$th $x$ derivative at $x_0$, $a\le s$,
either vanishes or has a factor $x_0^{3j-a+i}$ with nonnegative
exponent. Their coefficients use at most $s$ derivatives of $G$.

For the Taylor remainder, a chain-rule term using $i$ direct
$x$ derivatives and $k$ $\lambda$ derivatives is bounded by
\[
 C_s\|G\|_{C^s}
 x_0^{3(s-i-k)/2}x_0^{3k/2-(a-i)}
 =C_s\|G\|_{C^s}x_0^{3s/2-a-i/2}.
\]
Again the exponent is nonnegative. At exponent zero the remainder
tends to zero by continuity of the top Taylor derivatives. The
polynomial terms tend to the appropriate endpoint Taylor jet, so
these are bounds for continuous derivatives up to $x=0$.
Apply this to the even part of $G$ to prove the first estimate.
Tangential derivatives commute with the construction and reduce
the remaining derivative budget by their order.

For the second estimate the smooth even function
\[
 K(x,\lambda,\theta)
 =\frac12\int_{-1}^1\partial_\lambda G(x,t\lambda,\theta)\,dt
\]
satisfies $\|K\|_{C^s}\le C_s\|G\|_{C^{s+1}}$ and
$\mathcal O G=K(x,x^{3/2},\theta)$. Apply the even estimate to
$K$. Both operations are linear in $G$, so the same bounds commute
with parameter differentiation.
\end{proof}

Applying the lemma to the common flow gives
\[
 \sigma=x^{3/2}\chi_0(x,\theta),\quad
 \chi_0(0,\theta)=\tfrac23c_0(0,\theta)h_L'(\theta)>0,\quad
 \xi_L=x(3\chi_0/2)^{2/3}.
\]
Thus the root coordinate is smooth with nonzero radial derivative.
The right root uses its common boundary chart $h_R$ in the same way.
Rescaling from the unit rectangle to a root collar introduces powers
of its length and reciprocal, all bounded by $\mathcal E_m^{C_s}$.
Their complete product costs one $M$ order. The even substitution
uses $s$ input derivatives and the odd divided difference uses
$s+1$; these costs remain fixed as the output order increases.

\begin{figure}[htbp]
\centering
\begin{tikzpicture}[x=2.5cm,y=2.0cm,font=\small,>=Stealth]
\draw[->] (0,0)--(1.45,0) node[right] {$x$};
\draw[->] (0,0)--(0,1.45);
\draw[thick,blue!65!black] plot[domain=0:1.2,samples=50] (\x,{pow(\x,1.5)});
\draw[thick,orange!85!black] (0,0)--(1.2,1.2);
\node[blue!65!black,anchor=west] at (.74,.48) {$\sigma\sim x^{3/2}$};
\node[orange!85!black,anchor=east] at (.72,1.1) {$\xi\sim x$};
\draw[dashed] (1,0)--(1,1);
\node[align=left,anchor=west] at (1.65,.85)
 {The travel angle has zero derivative.\\
  Its $2/3$ power has a nonzero derivative.};
\end{tikzpicture}
\caption{Regularization at a turning circle. The graph displays the
leading powers with unit coefficients. The positive coefficient of
$x$ in $\xi$ is controlled uniformly at every turning circle.}
\label{fig:root-regularization}
\end{figure}

Write $\Psi_i=(\xi_i,\varphi)$ and
$\alpha_i=(\delta\circ\Psi_i^{-1})/\sqrt{\xi_i}$.
The factor is smooth and positive, also at the root, since
\[
 \alpha_i=\frac{a|\det D\Psi_i|}{U_\theta V_\theta}.
\]
For $u=(h_iW_i)\circ\Psi_i$, $h_i=\alpha_i^{-1/2}$, the equation is
\begin{equation}\label{eq:fixed-strip-airy-equation}
 \begin{split}
 (\partial_{\xi_i}^2-\xi_i\partial_\varphi^2+V_i)W_i&=F_i,\\
 V_i&=-\alpha_i^{-1/2}
       (\partial_{\xi_i}^2-\xi_i\partial_\varphi^2)\alpha_i^{1/2},\\
 F_i&=h_i((f/|\det D\Psi_i|)\circ\Psi_i^{-1}).
 \end{split}
\end{equation}
Indeed the density in these coordinates is
$\alpha_i\operatorname{diag}(1,-\xi_i)$; the positive conjugation
cancels its first derivatives. The source formula preserves
$\int uf\,dr\,d\theta=\int W_iF_i\,d\xi_i\,d\varphi$.

The middle value is $\xi_c=(3T/4)^{2/3}$. There
$\alpha_L=\alpha_R$ and
$\partial_\sigma\log\alpha_R-\partial_\sigma\log\alpha_L
=\xi_c^{-3/2}$. Consequently value and conormal matching become
\begin{equation}\label{eq:fixed-strip-cut}
 \binom{W_R}{(W_R)_{\xi_R}}
 =\begin{pmatrix}1&0\\-1/(2\xi_c)&-1\end{pmatrix}
       \binom{W_L}{(W_L)_{\xi_L}}.
\end{equation}
This matrix is independent of the state. At a type root, with boundary
chart $h_{\mathrm{bd}}$, the physical pair is
\begin{equation}\label{eq:fixed-strip-root-pair}
 u=(h_iW_i)(0,h_{\mathrm{bd}}),\qquad
 p_{\mathrm{in}}=\frac{h_{\mathrm{bd}}'}{h_i(0,h_{\mathrm{bd}})}
       (\partial_{\xi_i}W_i+\partial_{\xi_i}\log h_i\,W_i)
                                             (0,h_{\mathrm{bd}}).
\end{equation}
The tangential shear cancels because $a\varphi_r+b\varphi_\theta=0$
at the root. This identity keeps the zero mode and the full Green
pairing, not just the principal characteristics.

For completeness, all reference coordinate reciprocals used here have
the promised bounds. At a root
\[
 \xi=x\left[\tfrac32\int_0^1\sqrt t\,c_i(xt)\,dt\right]^{2/3}.
\]
Outside its collar the Morse potential has a lower bound $cA_n^{24}$;
the upper and lower bounds for $v$ then bound
$|\xi_r|=v_0/\sqrt\xi$ above and below by powers of $A_n$.
Also $\xi_c\asymp A_n^{1/3}$. Differentiating these formulas and
the inverse identity gives $\mathcal E_n^{C_s}$ bounds for charts,
inverses, $\alpha_i=w_*|\xi_r|$, $h_i$, potentials and source
factors. The perturbed bounds follow from the flow and inverse
argument in Proposition~\ref{prop:uniform-physical-geometry}.

\subsection{Crossing the positive components}

On a positive component continue the signed root coordinate and put
$y_i=-\xi_i$. Build a radial reference coordinate $y\in[0,1]$
equal to $y_L$ near the left root and $1-y_R$ near the right.
This can be done explicitly by choosing positive $y_r$ with integral
one: use the prescribed derivatives on disjoint root collars and a
positive constant on their complement, with cutoff interpolation.
The collars have total prescribed coordinate increment less than
$1/4$, so that constant and its reciprocal are bounded by
$\mathcal E_m^C$. Interpolate the logarithms of the two root height
factors, choosing height factor one in the middle. All reference
coordinate, height and source factors retain the same jet bounds.

For a perturbed state extend the \emph{differences} of the negative
coordinate jets through degree four and of the logarithmic height
jets through degree three, using cutoff Taylor polynomials. This
leaves the reference map unchanged on a compact interval $J_*$
containing the second-mode normalization and metric control.
Relative closeness of the differential gives a global diffeomorphism:
it is proper, locally invertible, and of boundary degree one.
The map and its inverse have the estimates already proved.

Writing the positive height factor in physical coordinates as
$h_{\mathbf x}$ and its chart as $\Psi_{\mathbf x}$, the positive
canonical equation for $-\widehat Lu=F_{\rm ph}$ is
\begin{equation}\label{eq:pp-fixed-canonical-source}
 P_{\mathbf x}U
 =-\left(\frac{h_{\mathbf x}}{|J_{\mathbf x}|}
       \widehat L[h_{\mathbf x}(U\circ\Psi_{\mathbf x})]\right)
             \circ\Psi_{\mathbf x}^{-1},\qquad
 F=\left(\frac{h_{\mathbf x}F_{\rm ph}}{|J_{\mathbf x}|}\right)
             \circ\Psi_{\mathbf x}^{-1}.
\end{equation}
These formulas again preserve the pairing. The continued negative
principal part is $-\partial_{y_i}^2-y_i\partial_\varphi^2$.
The specified jets match principal densities through order two.
Thus, with $d=y(1-y)$,
\begin{equation}\label{eq:readable-positive-difference}
 P_{\mathbf x}-P_*=-\partial_a(d^3e_{\mathbf x}^{ab}\partial_b)
                           +\delta V_{\mathbf x}.
\end{equation}
There is no boundary condition on the potential. Division by $d^3$
has a fixed cost: if $F,F',F''$ vanish at zero, then
\[
 F(y)/y^3=\tfrac12\int_0^1(1-t)^2F'''(ty)\,dt.
\]
Each further derivative differentiates the integrand; the $L^2$
rescaling factor $t^{-1/2}$ is integrable. The right end is identical.
This verifies the final divided-coefficient assertion of the uniform
geometric calculus. All operations use at most two height derivatives,
one root divided difference, four boundary jets, three divisions and
fixed trace and embedding costs; their sum fixes $d_*$.

The reference positive operator is
\[
 P_*=-\partial_y(a\partial_y)-c\partial_\varphi^2+V,
 \qquad a=h_*^2w_*y_r,\quad c=\frac{h_*^2w_*'}{ry_r}.
\]
It obeys $a\ge\mathcal E_m^{-C}$ and
$c\ge\mathcal E_m^{-C}d$, with all fixed coefficient jets bounded
by $\mathcal E_m^{C_s}$. Near a root it is exactly $a=1,c=y_i$;
on the rest this follows from the positive-model and transition lower
bounds for $w_*'/w_*$. These are the coefficients of the next energy
estimate.

At a type root the adjacent charts and height factors have the same
value and first physical derivative. Formula~\eqref{eq:fixed-strip-root-pair}
therefore becomes equality of canonical values and opposite inward
derivatives. Together with \eqref{eq:fixed-strip-cut} this is one
fixed transmission domain. Its smooth members really are smooth in
the original variables: the value matches tangential derivatives,
the conormal matches the first normal derivative, and the equation
solves for $a u_{rr}$ since $a>0$. Repeated normal differentiation
then matches every higher jet. This observation is what permits us
to use a common closed operator domain in the next section.
